\documentclass[11pt,a4paper]{article}

\usepackage[T1]{fontenc}
\usepackage[utf8]{inputenc}
\usepackage[french]{babel}
\usepackage{lmodern}
\usepackage{microtype}
\usepackage[a4paper,margin=2.4cm]{geometry}
\usepackage{xcolor}
\usepackage{booktabs}
\usepackage{tabularx}
\usepackage{hyperref}
\usepackage[backend=biber,style=authoryear,sorting=nyt,maxbibnames=99]{biblatex}

\definecolor{esilvblue}{HTML}{153A5B}
\definecolor{esilvcyan}{HTML}{1C7897}
\hypersetup{colorlinks=true,linkcolor=esilvblue,urlcolor=esilvcyan,citecolor=esilvblue}
\addbibresource{references.bib}

\title{\textcolor{esilvblue}{Projet PI\textsuperscript{2} n\textsuperscript{o} 176}\\
\Large Dynamic of RND and Forecasting the Price of Derivatives\\[0.4em]
\large Bibliographie structurée}
\author{Rayan Ait Arkoub, Matthieu Balliste, Ilan Chahdi,\\
Nathan Denoux et Jordi Paris}
\date{14 octobre 2026}

\begin{document}
\maketitle

\begin{abstract}
Cette bibliographie couvre les cinq blocs du projet : constitution des données,
estimation des distributions risque-neutre, mesure de leur stabilité, modélisation
de leur dynamique et application à une stratégie de couverture ou de trading.
Elle reprend les références de la fiche projet et les complète par les documents
fournis par l'équipe et quelques sources directement utiles à l'implémentation.
\end{abstract}

\section*{Organisation des lectures}

\begin{tabularx}{\textwidth}{@{}p{0.22\textwidth}X@{}}
\toprule
\textbf{Bloc} & \textbf{Références prioritaires} \\
\midrule
Fondements et extraction & Breeden et Litzenberger (1978) ; Figlewski (2010, 2018) \\
Lissage et robustesse & Shimko (1993) ; Taboga (2016) ; BSIC (2026) \\
Méthodes alternatives & Cui et Yu (2021) ; Cui, Liu et Yao (2025) \\
Entropie et calibration & Cover et Thomas (1991) ; Avellaneda et al. (2001) \\
Dynamique et prévision & Giacomini, Härdle et Krätschmer (2009) \\
Application économique & Martin (2017) ; méthodologie VIX de Cboe \\
\bottomrule
\end{tabularx}

\section*{Priorités de lecture}

\begin{enumerate}
  \item Figlewski (2018), pour cartographier l'ensemble du champ.
  \item Breeden et Litzenberger (1978), pour le résultat théorique central.
  \item BSIC (2026), Figlewski (2010) et Shimko (1993), pour le pipeline empirique.
  \item Cui et Yu (2021), pour la comparaison avec une approche Fourier-cosinus.
  \item Cover et Thomas (1991), puis Giacomini et al. (2009), pour la stabilité et la prévision.
  \item Taboga (2016), Avellaneda et al. (2001) et Martin (2017), pour l'incertitude et l'application.
\end{enumerate}

\section*{Point méthodologique central}

Pour une option européenne et un taux déterministe, l'identité de
Breeden--Litzenberger relie la densité risque-neutre au prix du call :
\[
q_t(K,T)=e^{r(T-t)}\frac{\partial^2 C_t(K,T)}{\partial K^2}.
\]
La double dérivation amplifie le bruit des cotations. L'interpolation, le lissage
sans arbitrage, le traitement des queues et les contrôles de cohérence probabiliste
doivent donc faire partie intégrante du protocole expérimental.

\nocite{*}
\printbibliography[title={Références}]

\end{document}
